\documentclass[11pt]{article}

\usepackage[margin=1in]{geometry}
\usepackage{amsmath,amssymb,amsthm,mathtools}
\usepackage{microtype}
\usepackage{enumitem}
\usepackage[hidelinks]{hyperref}
\usepackage{xurl}
\usepackage{booktabs}

\newtheorem{theorem}{Theorem}
\newtheorem{lemma}{Lemma}
\newtheorem{proposition}{Proposition}
\newtheorem{definition}{Definition}
\newtheorem{remark}{Remark}

\DeclareMathOperator{\TV}{TV}

\title{Receipt Redaction Bounds:\\ Contraction Coefficients for Publishable Privacy Artifacts}
\author{}
\date{Draft v0.10 (2026-02-23)}

\begin{document}
\maketitle

\begin{abstract}
Several papers in this series treat anonymity as a \emph{publishable artifact}: a receipt, certificate, or transcript that third parties can verify.
Publishing such artifacts creates a new channel: the receipt itself can leak information unless it is carefully redacted.

We give a compact theory module for bounding that incremental leak.
We model a redaction as a channel (Markov kernel) from sensitive internal state to published fields and bound the induced total-variation (TV) distinguishability in terms of a \emph{contraction coefficient}.
We highlight a practical corollary: if a published field is produced by an $\varepsilon$-local-DP (LDP) mechanism, then its TV contribution is at most $\tanh(\varepsilon/2)$, independent of the underlying state distribution.
This yields a clean, composable knob for ``how much you can publish'' when attaching receipts to releases, audits, and monitoring logs.
\end{abstract}

\paragraph{Scope.}
We give small, reusable lemmas about \emph{receipt redaction}: how much information can be removed from a public release ledger while still preserving the ability to verify key claims.

\paragraph{Imports.}
We import the notion of ``release receipts'' and their use as public commitments (Release~A), and the broader ABOM/OINL interface contract (Anonymity~C).\cite{series:releaseproperty,series:abomoinl}

\paragraph{Exports.}
Exports two redaction bounds: (i) conditions under which removing a field preserves the ability to recompute a verifier hash, and (ii) a compositional rule for combining redactions across nested receipt layers (e.g., per-build receipts inside a transparency log).

\paragraph{Artifact policy.}
This archive is source-only; supporting artifacts are available on request as a bundle committed by ABOM/OINL and governed by the series artifact policy (Synthesis~6).\cite{series:artifactpolicy}

\section{Receipt fields as a channel}
Let $X$ denote a sensitive internal object (e.g., a raw trace, a routing transcript, a congestion time series, or a labeled audit episode),
and let $R$ denote a published receipt field (or a tuple of fields).
A \emph{redaction} is a (possibly randomized) mapping from $X$ to $R$.

\begin{definition}[Receipt channel and Dobrushin coefficient]
A receipt channel is a Markov kernel $K(\cdot\mid x)$ from $\mathcal{X}$ to $\mathcal{R}$.
Its (TV) contraction coefficient is
\[
\eta(K) \;:=\; \sup_{x,x'} \TV\!\left(K(\cdot\mid x),K(\cdot\mid x')\right)\in[0,1].
\]
\end{definition}

\begin{lemma}[TV cannot increase through a channel]
For any two distributions $P,Q$ on $\mathcal{X}$ and any channel $K$,
\[
\TV(PK,QK) \le \eta(K)\,\TV(P,Q).
\]
In particular, $\TV(PK,QK)\le \TV(P,Q)$.
\end{lemma}

\begin{proof}
This is a standard Dobrushin-contraction bound.
One proof uses the dual characterization of TV as a supremum over tests $\varphi:\mathcal{R}\to[0,1]$ and the induced test $\psi(x):=\mathbb{E}[\varphi(R)\mid X=x]$.
Then $|\mathbb{E}_{PK}[\varphi]-\mathbb{E}_{QK}[\varphi]|=|\mathbb{E}_{P}[\psi]-\mathbb{E}_{Q}[\psi]| \le \TV(P,Q)\,\|\psi\|_{\mathrm{Lip}}$,
and $\|\psi\|_{\mathrm{Lip}}\le \eta(K)$ by definition of $\eta(K)$.
\end{proof}

\paragraph{Interpretation.}
If you publish a receipt field $R=g(X)$ with \emph{no randomization}, then worst-case $\eta(K)=1$ and the field can leak as much as the raw state.
To make $\eta(K)$ small, the release must add noise or otherwise ensure that different internal states induce similar laws for the published object.

\section{Local DP implies a clean TV cap}
Local differential privacy is a natural way to randomize a published receipt field: it is an \emph{interface-level} promise that does not depend on the deployment distribution.

\begin{definition}[$\varepsilon$-local DP]
A channel $K$ is $\varepsilon$-local-DP if for all $x,x'$ and measurable $A\subseteq\mathcal{R}$,
\[
K(A\mid x) \le e^{\varepsilon} K(A\mid x').
\]
\end{definition}

\begin{proposition}[TV cap for $\varepsilon$-LDP channels]\label{prop:tv-ldp}
If $K$ is $\varepsilon$-local-DP, then
\[
\eta(K)\le \frac{e^{\varepsilon}-1}{e^{\varepsilon}+1}=\tanh(\varepsilon/2).
\]
\end{proposition}

\begin{proof}
Fix $x,x'$ and let $P=K(\cdot\mid x)$, $Q=K(\cdot\mid x')$.
By $\varepsilon$-LDP, the likelihood ratio satisfies $dP/dQ \in [e^{-\varepsilon},e^{\varepsilon}]$ wherever densities exist (and the analogous statement holds in the purely discrete case).
The worst-case TV under a bounded likelihood ratio is attained by a two-point construction and equals $(\alpha-1)/(\alpha+1)$ where $\alpha=\sup dP/dQ$.
Setting $\alpha=e^{\varepsilon}$ yields the claim.
\end{proof}

\begin{remark}[Why this is useful]
Many papers in this series already bound leakage from the \emph{protocol} channel.
Proposition~\ref{prop:tv-ldp} bounds leakage from the \emph{published certificate/receipt} channel itself.
Operationally: to publish an artifact field, pick an $\varepsilon$ for that field and record it in the receipt; then the field contributes at most $\tanh(\varepsilon/2)$ TV in the worst case.
\end{remark}

\section{Composing multiple receipt fields}
Receipts are typically tuples: $R=(R_1,\dots,R_k)$.
Let $K_i$ be the channel for field $i$.

\begin{proposition}[Union-style bound]\label{prop:union}
If the receipt fields are generated conditionally independently given $X$ (i.e., the channel factorizes as $K=\prod_i K_i$), then
\[
\eta(K) \le 1-\prod_{i=1}^k (1-\eta(K_i)) \le \sum_{i=1}^k \eta(K_i).
\]
\end{proposition}

\begin{proof}
For fixed $x,x'$, couple each field channel so that $R_i=R_i'$ with probability at least $1-\TV(K_i(\cdot\mid x),K_i(\cdot\mid x'))$.
Under conditional independence, the product coupling makes all fields equal with probability at least $\prod_i (1-\TV_i)$.
Thus the TV of the product laws is at most the probability that at least one field differs, giving the first bound.
The second is $1-\prod_i (1-a_i)\le \sum_i a_i$.
\end{proof}

\paragraph{Practical guideline.}
Keep receipt fields \emph{few} and \emph{separable}.
If a single derived field depends on many raw internals, its effective contraction coefficient can be close to 1 unless you explicitly add noise.

\section{Examples for publishable anonymity artifacts}
\paragraph{(1) Publishing bucketed tail summaries.}
Suppose a system wants to publish a tail indicator (e.g., whether a latency exceeded a threshold) for monitoring.
Publishing a raw Bernoulli outcome can leak heavily if it is correlated with the secret.
Applying randomized response with LDP parameter $\varepsilon$ yields an explicit $\tanh(\varepsilon/2)$ TV cap for that field.

\paragraph{(2) Publishing histogram sketches.}
If a paper wants to publish a small sketch of an empirical distribution (e.g., a congestion witness histogram or a destination-trace bucket count),
then using an LDP mechanism per bin (or a private count sketch) yields an explicit per-field TV bound.
Proposition~\ref{prop:union} then gives a simple end-to-end cap.

\paragraph{(3) Receipt design in the series.}
The release-receipt paper should record \emph{receipt-field privacy knobs} (e.g., LDP parameters or explicit redaction policies) as first-class items in the exported receipt.
The closed-view auditing paper should treat published canary and monitoring summaries as receipt fields and budget them similarly.

\section{What this note does \emph{not} do}
This note does not choose an optimal redaction for any specific system.
It provides a small, reusable interface: if you can bound (or choose) a contraction coefficient for a published artifact field, you can budget the receipt channel alongside the protocol channel.

\begin{thebibliography}{99}\small

\bibitem{series:artifactpolicy}
Anonymous.
\newblock Artifact and Disclosure Policy for the One-True-Archive (Source-Only Public Release).
\newblock Series synthesis note (Synthesis~6), 2026. (This archive.)

\bibitem{series:releaseproperty}
Anonymous.
\newblock Release Property Ledger Receipts (Commitments, Redaction, and Verifiability).
\newblock Release and Destination series Paper~A, 2026. (This archive.)

\bibitem{series:abomoinl}
Anonymous.
\newblock ABOM and OINL for Anonymity Claims.
\newblock Anonymity series Paper~C, 2026. (This archive.)

\end{thebibliography}

\end{document}
